\documentclass[12pt,a4paper]{amsart}

\usepackage{amsmath,amssymb,amsthm}
\usepackage{mathtools}
\usepackage{hyperref}
\usepackage{geometry}
\geometry{margin=2.5cm}

% -------------------------------------------------------
\newtheorem{theorem}{Theorem}[section]
\newtheorem{lemma}[theorem]{Lemma}
\newtheorem{corollary}[theorem]{Corollary}
\newtheorem{proposition}[theorem]{Proposition}
\theoremstyle{definition}
\newtheorem{definition}[theorem]{Definition}
\newtheorem{remark}[theorem]{Remark}

\newcommand{\Z}{\mathbb{Z}}
\newcommand{\euler}{\varphi}

% -------------------------------------------------------
\begin{document}
% -------------------------------------------------------

\title{Every Solution to Lehmer's Totient Problem is Squarefree:\\ A Short Proof}

\author{Kurt Ingwer}
\address{Specialist Mathematics Project, Version Build 44}
\email{specialist-maths@localhost}
\date{\today}

\begin{abstract}
Lehmer's totient problem (1932) asks whether the condition $\euler(n) \mid (n-1)$
forces $n$ to be prime.  A composite solution would be called a \emph{Lehmer number}.
We give a short, elementary proof that every Lehmer number must be squarefree.
The argument rests on two steps: if $p^2 \mid n$, then $p \mid \euler(n)$;
combined with $\euler(n) \mid (n-1)$, this gives $p \mid (n-1)$; but also $p \mid n$,
so $p \mid 1$, which is impossible.  This result is fundamental for the study of
Lehmer's problem and implies that any potential counterexample is a product of distinct primes.
\end{abstract}

\maketitle

% -------------------------------------------------------
\section{Introduction}
% -------------------------------------------------------

In 1932, Lehmer \cite{Lehmer1932} asked whether the congruence
$\euler(n) \mid (n-1)$ (where $\euler$ is Euler's totient function) can hold
for a composite integer~$n$.  Primes always satisfy $\euler(p) = p-1 \mid p-1$;
the problem is whether this characterises primes.

\begin{definition}
A composite integer~$n$ with $\euler(n) \mid (n-1)$ is called a
\emph{Lehmer number} or \emph{Lehmer pseudoprime}.
\end{definition}

No Lehmer number has ever been found.  Cohen and Hagis \cite{Cohen1980} proved that
any such number must have at least~$13$ prime factors; Banks and Luca \cite{Banks2009}
improved this to~$15$.

The present note establishes the most basic structural property:
squarefreeness.

% -------------------------------------------------------
\section{Main Result}
% -------------------------------------------------------

\begin{theorem}
\label{thm:sqfree}
Every Lehmer number is squarefree.
\end{theorem}

\begin{proof}
Suppose $n$ is a Lehmer number and $p^2 \mid n$ for some prime~$p$.

Since $p^2 \mid n$, we have $p \mid \euler(n)$.
(Recall: if $p^a \| n$ with $a \ge 2$, then $p^{a-1}(p-1) \mid \euler(n)$,
so in particular $p \mid \euler(n)$; even for $a = 2$ this gives $p(p-1) \mid \euler(n)$.)

From $\euler(n) \mid (n-1)$ and $p \mid \euler(n)$:
\[
  p \mid \euler(n) \mid (n-1), \quad \text{so} \quad p \mid (n-1).
\]
But also $p \mid n$ (since $p^2 \mid n$), hence
\[
  p \mid n - (n-1) = 1.
\]
This is impossible for $p \ge 2$ — a contradiction.
Therefore no prime~$p$ satisfies $p^2 \mid n$, i.e.\ $n$ is squarefree.
\end{proof}

\begin{remark}
The same argument applies to the more general condition $M \cdot \euler(n) = n - 1$
for any fixed positive integer~$M$ (a generalisation studied by Hagis \cite{Hagis1988}).
Squarefreeness follows identically.
\end{remark}

% -------------------------------------------------------
\section{Consequences}
% -------------------------------------------------------

\begin{corollary}
Every Lehmer number has the form $n = p_1 p_2 \cdots p_k$ with distinct primes
$p_1 < p_2 < \cdots < p_k$ and $k \ge 2$.
\end{corollary}

\begin{corollary}
For a Lehmer number $n = p_1 \cdots p_k$, the totient is
$\euler(n) = (p_1-1)(p_2-1)\cdots(p_k-1)$, and the Lehmer condition becomes
\[
  (p_1-1)(p_2-1)\cdots(p_k-1) \mid p_1 p_2 \cdots p_k - 1.
\]
\end{corollary}

This explicit form of the condition is the basis for all further analysis of
Lehmer numbers, including the proofs that no semiprime \cite{Ingwer2026b} and
no three-prime number \cite{Ingwer2026lehmer3} can be a Lehmer number.

% -------------------------------------------------------
\section{Comparison with the Giuga Case}
% -------------------------------------------------------

The parallel structure with Giuga pseudoprimes is striking.  For both problems:
\begin{itemize}
  \item Squarefreeness is proved by showing $p^2 \mid n$ leads to $p \mid 1$.
  \item For Giuga: $p^2 \mid n \Rightarrow p \mid \tfrac{n}{p} \Rightarrow \tfrac{n}{p}-1 \equiv -1 \pmod{p} \Rightarrow p \mid -1$.
  \item For Lehmer: $p^2 \mid n \Rightarrow p \mid \euler(n) \Rightarrow p \mid (n-1) \Rightarrow p \mid 1$.
\end{itemize}
Both use the fundamental incompatibility between $p \mid n$ and $p \mid (n-1)$.

% -------------------------------------------------------
\begin{thebibliography}{10}

\bibitem{Lehmer1932}
D.~H.~Lehmer,
\emph{On Euler's totient function},
Bull.\ Amer.\ Math.\ Soc.\ \textbf{38} (1932), 745--751.

\bibitem{Cohen1980}
G.~L.~Cohen and P.~Hagis,
\emph{On the number of prime factors of $n$ if $\phi(n) \mid (n-1)$},
Nieuw Arch.\ Wisk.\ (3) \textbf{28} (1980), 177--185.

\bibitem{Hagis1988}
P.~Hagis,
\emph{On the equation $M \cdot \phi(n) = n - 1$},
Nieuw Arch.\ Wisk.\ (4) \textbf{6} (1988), 255--261.

\bibitem{Banks2009}
W.~D.~Banks and F.~Luca,
\emph{Composite integers $n$ for which $\phi(n) \mid n - 1$},
Acta Math.\ Sin.\ (Engl.\ Ser.) \textbf{25} (2009), 1703--1710.

\bibitem{Ingwer2026b}
K.~Ingwer,
\emph{No semiprime satisfies Lehmer's totient condition},
Specialist Mathematics Project (2026).

\bibitem{Ingwer2026lehmer3}
K.~Ingwer,
\emph{Lehmer's totient problem for products of three primes},
Specialist Mathematics Project (2026).

\end{thebibliography}

\end{document}
